\section{Knowledge Circuits Discovery in Transformers}
\subsection{Knowledge Circuits Construction}
Unlike previous work~\cite{dai-etal-2022-knowledge,meng2022locating}, which managed to find out the specific areas that store knowledge, we pay extra heed to the information flow that activates subsequent knowledge for answering questions.
Similar to \cite{goldowsky2023localizing, wang2023interpretability}, we write language model as a graph consisting of the input, the output, attention heads, and MLPs by considering a ``residual rewrite'' of the model's computational structure.
For example, this residual rewrite gives us a nearly-dense graph in GPT2-medium: one between every pair of (attention head, MLP, input, and output) nodes, except for attention heads in the same layer, which do not communicate with each other.
% $A(\mathbf{r})=\left(\operatorname{softmax}\left(\mathbf{r}^{\top} W_Q^{\top} W_K \mathbf{r}\right) \otimes W_O W_V\right) \cdot \mathbf{r}$
% We can write any model as a connected directed acyclic graph (DAG) with source nodes representing the model’s (typically vector-valued) input, sink nodes representing the model’s output, and intermediate nodes representing units of computation.
In our paper, we concentrate on the task of answering factual open-domain questions, where the goal is to predict a target entity $o$ given a subject-relation pair $(s, r)$.
A knowledge triplet $k=(s, r, o)$ is often presented to the model in the form of a natural language prompt for next token prediction (e.g.,  \nl{The official language of France is \textunderscore\textunderscore\textunderscore\textunderscore}). 
The model $\mathcal{G}$ is expected to generate the target entity, which is consistent with the language model's pretraining format.
To identify the circuit that is critical for predicting the target entity $o$ for a given subject-relation pair $(s, r)$, we ablate each special edge $e_{i} = (n_{x},n_{y})$ in the computation graph $\mathcal{G}$. 
We then measure the impact of ablating the edge (\textbf{zero ablation} in our implementation) on the model's performance using the MatchNLL loss \cite{acdc} for the target $o$:
\begin{equation}
    S(e_{i}) =  \log(\mathcal{G}(o|(s,r)))- \log(\mathcal{G}/e_{i}(o|(s,r)))
    \label{metric_S}
\end{equation}
If the score $S(e_{i})$ is less than the predefined threshold $\tau$, we consider the edge to be non-critical and remove it from the computation graph, updating the temporary circuit $\mathcal{C}_{temp} \leftarrow\mathcal{G}/e_{i}$. 
We first sort the graph by topological rank following \citet{acdc} and  traverse all edges in this manner, We derive a circuit $\mathcal{C}_{k}$ that contributes to representing the knowledge necessary to answer the factual question:
\begin{equation}
    \mathcal{C}_{k} = <N_{k}, E_{k}>
\end{equation}
Here, $\mathcal{C}_{k}$ is the circuit for the knowledge triplet $k$, consisting of the nodes $N_{k}$ and edges $E_{k}$ that are essential for predicting the target entity $o$ given the subject-relation pair $(s, r)$.


\subsection{Knowledge Circuits Information Analysis}
Once we have identified the knowledge circuit, we delve deeper into the specific roles and behaviors of each node and edge within the computation graph. 
Our goal is to comprehend the processing and contribution of each node $n_{i}$ to the functionality of the circuit.
Drawing on the methodologies of previous studies \cite{yu2024locating,katz-belinkov-2023-visit,nostalgebraist_interpreting_nodate}, we begin by applying layer normalization to the output of each node $n_{i}$ and then map it into the embedding space. 
This is achieved by multiplying the layer-normalized output by the unembedding matrix ($\mathbf{W}_U$) of the language model:
$\mathbf{W}_U \operatorname{LN}(n_{i})$.
This transformation allows us to inspect how each component writes information to the circuit and how it influences subsequent computational steps. 
% Additionally, we measure the difference between the input and output of the node when mapped into the embedding space to discern the type of information processing that occurs:$\mathbf{W}_U \operatorname{LN}(n_{i})$.
By understanding the nodes' behavior in the circuit, we can better comprehend the circuit's structure and the key points where information is aggregated and disseminated.
% $\Delta n_{i} = \mathbf{W}_U \operatorname{LN}(n_{i}) - \mathbf{W}_U \operatorname{LN}(\text{Input}(n_{i}))$
% Furthermore, we analyze the in-degree and out-degree of the nodes within the circuit to identify the most critical nodes and to search for potential bottlenecks in the flow of information.

\begin{table}[ht]
\small
\vspace{-3mm}
    \centering
    \small
    % \renewcommand{\arraystretch}{1.3}
    \caption{\texttt{\textbf{Hit@10}} of the Original and Circuit Standalone performance of knowledge circuit in GPT2-Medium. \textbf{The result for $D_{val}$ being 1.0 indicates that we select the knowledge for which the model provides the correct answer to build the circuit.} } 
    \resizebox{\linewidth}{!}{
    \begin{tabular}{lccccccc}
    \toprule
    \multirow{2}{*}{\textbf{Type}}  & \multirow{2}{*}{\textbf{Knowledge}} & \multirow{2}{*}{\textbf{\#Edge}} &\multicolumn{2}{c}{$D_{val}$} &  \multicolumn{3}{c}{\texttt{\textbf{$D_{test}$}}} \\ \cmidrule{4-8}
      & && Original($\mathcal{G}$) &  Circuit($\mathcal{C}$) & Original($\mathcal{G}$) & Random & Circuit($\mathcal{C}$)  \\ \midrule
    \multirow{3}{*}{Linguistic}&   Adj Antonym & 573 & 0.80 & {  1.00 \textcolor{red}{$\uparrow$}} & 0.00 & 0.00 & {  0.40  \textcolor{red}{$\uparrow$}} \\
     & word first letter & 432 & 1.00 & 0.88 & 0.36 & 0.00 & 0.16 \\
     & word last letter & 230 & 1.00 & 0.72 & 0.76& 0.00 & 0.76 \\ \midrule
    \multirow{3}{*}{Commonsense} &  object superclass & 102 & 1.00 & 0.68 & 0.64& 0.00 & 0.52\\
    & fruit inside color & 433 & 1.00 & 0.20 & 0.93 & 0.00 & 0.13 \\
        & work location & 422 & 1.00 & 0.70 & 0.10 & 0.00& 0.10 \\ \midrule

    \multirow{4}{*}{Factual}   & Capital City  & 451 & 1.00 & 1.00 & 0.00 & 0.00& 0.00 \\
     &   Landmark country & 278 & 1.00 & 0.60 & 0.16 & 0.00 & {  0.36 \textcolor{red}{$\uparrow$}} \\
      & Country Language & 329 & 1.00 & 1.00 & 0.16 & 0.00 & {  0.75 \textcolor{red}{$\uparrow$}} \\
       & Person Native Language & 92 & 1.00 & 0.76 &  0.50  & 0.00 & {  0.76 \textcolor{red}{$\uparrow$}} \\ \midrule
       
     \multirow{4}{*}{Bias} &  name religion & 423 & 1.00 & 0.50 & 0.42 & 0.00 & 0.42 \\
    & occupation age & 413 & 1.00 & 1.00 & 1.00 & 0.00 & 1.00 \\
        & occupation gender & 226 & 1.00 & 0.66 & 1.00& 0.00 & 0.66 \\ 
     & name birthplace &  276 & 1.00 & 0.57 & 0.07& 0.00 & {  0.57 \textcolor{red}{$\uparrow$}} \\ \midrule   
       \textbf{Avg}  &  &  & 0.98 & 0.73 &  0.44 & 0.00 & {  0.47 \textcolor{red}{$\uparrow$}} \\
       \bottomrule
    \end{tabular}}
    \label{tab:completeness}
\end{table}

\subsection{Knowledge Circuits Experimental Settings}


% \paragraph{Faithfulness}
% Faithfulness refers to the consistency of the knowledge circuit across different expressions of the same knowledge. To assess faithfulness, we examine whether circuits obtained from different knowledge prompts $\mathcal{C}_{P1}$ and $\mathcal{C}_{P2}$ are similar. This is quantified using a similarity measure $S$:
% \begin{equation}
%     S(\mathcal{C}_{P1}, \mathcal{C}_{P2})
% \end{equation}
% A high degree of similarity between circuits derived from different prompts indicates a faithful representation of the underlying knowledge, as the circuit should capture the essential components of the knowledge regardless of how it is expressed.



\paragraph{Implementations.}
\label{subsec:exp_setting}
We conduct experiments on GPT-style models, including GPT-2 medium and large.
We also conduct primary experiments on TinyLLaMA \cite{zhang2024tinyllama} to validate the effectiveness of different architectures.
We utilize the Automated Circuit Discovery \cite{acdc} toolkit to build a circuit as an initiative of our analysis and leverage transformer lens \cite{nanda2022transformerlens} to further analyze the results.
Specifically, we simply employ the MatchNLL \cite{acdc} as the metric to detect the effect of the given node and edge and use \textbf{zero ablation} to knock out the specific computation node in the model's computation graph.

\paragraph{Metrics.}


A discovered knowledge circuit is deemed an accurate representation of a specific area within the transformer's knowledge storage, thus, it should be capable of representing the knowledge independently.
Following \cite{acdc}, we leverage the completeness of a circuit, which refers to its ability to independently reproduce the behavior or predictions of the full model for the relevant tasks. 
This property is assessed by examining whether the identified subgraph corresponds to the underlying algorithm implemented by the neural network. 
To evaluate completeness, we first construct the circuit using the validation data $D_{val}$ for a specific knowledge type and then test its performance on the test split $D_{test}$ in isolation. 
By doing so, we can observe any changes in performance compared to the original model. 
We use the \texttt{\textbf{Hit@10}} metric to measure the rank of the target entity $o$ among the top 10 predicted tokens:
\begin{equation}
\texttt{\textbf{Hit@10}}=\frac{1}{|V|} \sum_{i=1}^{|V|} \mathrm{I}\left(\operatorname{rank}_{o} \leq 10\right)
\end{equation}
Here, $|V|$ represents vocabulary size, and $\operatorname{rank}_{o}$ is the rank of the target entity $o$ in predictions.

% check 是否


\paragraph{Dataset.}
In this work, we focus on the \textbf{knowledge that is already stored in the language model}.
We utilize the dataset provided by LRE~\cite{hernandez2024linearity} and consider different kinds of knowledge, including linguistic, commonsense, fact, and bias.
We evaluate whether the knowledge is present in the language model's parameters under zero-shot settings using the \texttt{\textbf{Hit@10}} metric to sample knowledge from the validation set, which is used to construct the knowledge circuit.
The data statistics are in Appendix~\ref{app_sec:implem_details}. 

